\begin{table}[t]
\caption{Preguntas de investigación}\label{tab:rq}
\footnotesize
\begin{tabularx}{\textwidth}{@{}L{0.9cm}Y{1.0}@{}}
\toprule
ID & Pregunta de investigación \\
\midrule
RQ1 & ¿Qué estrategias y algoritmos de perturbación proactiva se han propuesto para dificultar el entrenamiento no autorizado de modelos discriminativos con datos visuales, cómo se caracterizan y en qué tareas, conjuntos de datos y arquitecturas se han evaluado? \\
RQ2 & ¿Qué contramedidas se han evaluado frente a estas protecciones y bajo qué conocimientos, capacidades, acceso y recursos del adversario? \\
RQ3 & ¿Qué grado de protección persiste y qué recuperación del rendimiento logra el adversario tras aplicar las contramedidas evaluadas, en los distintos dominios y configuraciones experimentales? \\
RQ4 & ¿Qué referencias experimentales y métricas se utilizan para evaluar la robustez y la reversión adversaria, y en qué condiciones permiten comparar los resultados entre estudios y dominios? \\
\botrule
\end{tabularx}
\end{table}
